\section{Introduction}

The {\em throughput} of a sequence of instructions---the number processor clock cycles taken to execute the sequence when looped in steady state---determines how fast those instructions can process data.
Accurately predicting the throughput of a basic block\footnote{In this work, we focus specifically on \emph{basic blocks}, sequences of instructions with no branches or jumps.} is an essential requirement in many systems, to be able to predict and optimize runtime performance.
For instance, constraint--based register allocation and instruction scheduling~\cite{comb-inst-scheduling} relies on accurate throughput estimations, as do learning--based techniques like genetic algorithm based register allocation~\cite{ga-register-allocation} and reinforcement learning based instruction scheduling~\cite{rl-inst-scheduling}.

The alternative -- measuring throughput on demand by executing the basic block -- is too expensive for most compilers and learning--based solutions.
In practice, most systems employ analytical models to predict throughput.
For instance, the LLVM compiler team~\cite{llvm} recently merged\footnote{ lists.llvm.org/pipermail/llvm-dev/2018-March/121490.html} a command-line tool, llvm-mca~\cite{llvm-mca}, that exposes a machine model for throughput estimation.
Intel has also released a closed-source machine code analyzer, IACA~\cite{iaca}, which relies on internal knowledge of Intel's processor design.
These models are typically an order of magnitude faster than measuring a basic block's throughput.
However, manually writing an accurate and complete model is tedious, error-prone, and exceedingly difficult without knowledge of the exact mechanisms of the processor.

In the hunt for {\em accuracy}, developers build complicated models which must make significant tradeoffs with the model's {\em portability} and {\em speed}.

\paragraph{\bf Accuracy.}
Modern x86-64 Complex Instruction Set Computer (CISC) processors contain many hardware optimizations that significantly complicate building accurate analytical models.
In order to implement an instruction set \emph{architecture} (ISA) like x86-64, processors actually implement an underlying \emph{micro}architecture, a physical implementation of the ISA specification.
Processors translate instructions from the ISA to instructions in the latent microarchitectural language (termed micro-ops), then execute those micro-ops.
The micro-ops may undergo optimizations such as
\textit{micro-op fusion}, in which micro-ops of different instructions may be combined together;
\textit{out-of-order execution}, in which instructions can be executed in any semantics--preserving order;
\textit{register renaming}, where false dependencies can be broken to enable more parallel execution;
and many more vendor-specific optimizations.
This makes the prediction problem highly complex and non-linear.

\paragraph{\bf Portability.}
While ISAs like x86-64 stay relatively stable, processor vendors release updated processor implementations with different microarchitectures every few years.
For example, Intel released the Haswell and Skylake microarchitectures in 2013 and 2015 respectively for the x86-64 instruction set.
Each microarchitecture of a processor family has its own quirks and intricacies.
Manually writing a throughput estimator to support different microarchitectures requires rewriting instruction tables, resource utilization charts, and modeling microarchitectural optimizations, all of which are tedious and error-prone.
This is complicated by the vast, incomplete, and incorrect documentation for many processors, where understanding of these behaviors has to be obtained by reverse-engineering the processor.
Ideally, the throughput estimator should be able to automatically capture such intricacies with minimal human intervention.

\paragraph{\bf Speed.} A throughput estimator also needs to be fast.
Compilers need to search through many code blocks before emitting the fastest version of a given instruction sequence.
Running the basic blocks to get the ground truth throughput requires sandboxing and many iterations of execution to arrive at a consistent steady-state throughput estimate, which is impractical for real-time systems.

\subsection{\tool{}: A Data Driven Approach}

In this paper we introduce {\it \tool{} (Instruction THroughput Estimator using MAchine Learning)}, which takes a novel data--driven approach to predicting throughput for a block of instructions, inspired by recent advances in Deep Neural Networks (DNNs).
\tool{} models the throughput estimation problem as a regression task and leverages a DNN to learn to predict throughput by using a large corpus of labeled data, mapping assembly sequences to real valued throughputs.
More concretely, \tool{} uses a a hierarchical multiscale RNN~\cite{hihi-hierarchical, chung-hierarchical, baraldi-hierarchical}, which generates an independent embedding for each instruction, then sequentially combines the instruction embeddings to predict throughput.

We show that \tool{}'s learned model is significantly more accurate than the analytical models, dropping the mean absolute percent error by more than 50\% across all benchmarks, while still delivering fast estimation speeds.

To generate high-quality predictions, \tool{} needs only training data and a specification of the ISA, including the specification of instructions and their explicit and implicit operands (for instance, the instruction \texttt{push rax} in x86-64 pushes the register \texttt{rax} on to the stack and also {\em implicitly} modifies the stack pointer register, \texttt{rsp}).
Unlike analytical models, \tool{} learns any salient microarchitectural details that contribute to throughput on its own, without any explicit specification or modeling.

In this paper, we present the following contributions:
\setlist{nolistsep}
\begin{itemize}[noitemsep]
\item{\bf Data--Driven Throughput Estimation.}
  We present \tool{}, the first system for data-driven basic block throughput estimation, using a hierarchical multiscale RNN.

\item{\bf Evaluation:}
  We demonstrate that \tool{} is more accurate than the state-of-the-art analytical throughput estimators, while still attaining fast estimation speeds.
  We also show that \tool{}'s design is portable across multiple processor microarchitectures.

\item{\bf DNN Architecture Exploration.}
  We demonstrate that the proposed hierarchical multiscale RNN outperforms many other choices, including other neural network architectures specifically tailored to mimic the dependency patterns of instructions within a basic block.

\item{\bf Open Source Implementation.}
We have open-sourced our implementation of \tool{} at \url{https://github.com/psg-mit/Ithemal} in the hope that performance engineers and compiler designers can use and improve upon our approach.

\end{itemize}

{\noindent}{\tool} demonstrates that future systems can leverage data--driven techniques to either augment or fully replace manually developed throughput estimators.
